% Generated draft: requires mathematical and source review.
\documentclass[11pt]{article}
\usepackage[margin=25mm]{geometry}
\usepackage{fontspec}
\setmainfont{Latin Modern Roman}
\usepackage{amsmath,amssymb,mathtools,unicode-math}
\setmathfont{Latin Modern Math}
\usepackage{enumitem,needspace,xurl,hyperref}
\setlength{\parindent}{0pt}
\setlength{\parskip}{5pt}
\setlength{\emergencystretch}{3em}
\newcommand{\N}{\mathbb{N}}
\newcommand{\Z}{\mathbb{Z}}
\newcommand{\Q}{\mathbb{Q}}
\newcommand{\R}{\mathbb{R}}
\newcommand{\C}{\mathbb{C}}
\newcommand{\F}{\mathbb{F}}
\newcommand{\A}{\mathbb{A}}
\newcommand{\PP}{\mathbb{P}}
\newcommand{\E}{\mathbb{E}}
\newcommand{\eps}{\varepsilon}
\newcommand{\dd}{\,\mathrm{d}}
\begin{document}
\section*{UnsolvedMath ID: 30002620}
\textbf{Source problem code:} catalogue entry
\subsection{Definitions and mathematical statement}
A real \(k\)-current on \(\R^d\) is a continuous linear functional on smooth compactly supported differential \(k\)-forms. Its mass is
\[
M(T)=\sup\bigl\{|T(\omega)|:\omega\in C_c^\infty(\R^d;\Lambda^k),\ \sup_x\|\omega(x)\|_*\leq 1\bigr\},
\]
where \(\|\cdot\|_*\) denotes the comass norm. For \(k\geq1\), its boundary is defined by \(\partial T(\eta)=T(\mathrm d\eta)\). For \(k=0\), use the convention \(\partial T=0\). A current is normal if both \(M(T)\) and \(M(\partial T)\) are finite; write \(N_k(\R^d)\) for this space. An integral current is an integer-multiplicity rectifiable current of finite mass whose boundary is also integer-multiplicity rectifiable and has finite mass; write \(I_k(\R^d)\) for these currents.

\textbf{Decomposition question.} For every normal current \(T\in N_k(\R^d)\), does there exist a positive measure space \((L,\mathcal A,\lambda)\) and a family \((R_\ell)_{\ell\in L}\) of integral currents in \(I_k(\R^d)\) such that
\[
T=\int_L R_\ell\,\mathrm d\lambda(\ell),\qquad
M(T)=\int_L M(R_\ell)\,\mathrm d\lambda(\ell),\qquad
M(\partial T)=\int_L M(\partial R_\ell)\,\mathrm d\lambda(\ell)?
\]
Here the current-valued integral is interpreted weakly: for every smooth compactly supported \(k\)-form \(\omega\), the function \(\ell\mapsto R_\ell(\omega)\) is measurable and
\[
T(\omega)=\int_L R_\ell(\omega)\,\mathrm d\lambda(\ell).
\]
The two mass functions are also required to be measurable. The stated mass identity makes each scalar integral above absolutely convergent.

\textbf{Unresolved formulation qualifications.} The source does not specify the dimension range, coefficient convention, or meaning of the current-valued integral. The definitions above give the conventional real-current, integer-integral-current, weak-integral interpretation. The usual Euclidean dimension range is \(d\geq1\) and \(0\leq k\leq d\), but whether the source intends all these dimensions, including the endpoints, remains unspecified. No additional restriction on dimensions or on \((L,\mathcal A,\lambda)\) is supplied.

The supplied background describes a particular horizontal absolute-continuity question as answered positively by a later tangent-current theorem, while describing Morgan's broader normal-current decomposition and stronger distribution questions as not fully resolved. Those auxiliary questions and their hypotheses are not stated in the record, so no precise auxiliary assertion or implication for the displayed decomposition question can be extracted here. This is a statement-only draft requiring human mathematical review.
\subsection{Short English statement}
The question asks whether every normal current can be expressed as a measured family of integral currents while preserving both total mass and boundary mass exactly. The source leaves dimension and integration conventions unspecified. Its supplied dated assessment labels the topic partially solved, but does not supply enough detail to identify which cases of this decomposition question are settled.
\subsection{Sources}
\begin{itemize}
\item Supplied local UnsolvedMath registry record, permanent ID \(30002620\), source problem code \texttt{catalogue entry}, title \emph{Mass-Preserving Decomposition of Normal Currents}; registry location \texttt{data/unsolvedmath/problems.json}. This draft uses the supplied complete record, without independently inspecting the registry.
\item Supplied source citation: \emph{Calculus of Variations (2015)}, Oberwolfach Reports, pp.~1801--1860, DOI \texttt{10.4171/owr/2014/33}: \url{https://doi.org/10.4171/owr/2014/33}. The record describes Oberwolfach Reports as documenting mathematical workshops and their open-problem discussions, and records the proposed year as 2014.
\item Supplied literature reference: Alberti--Massaccesi--Stepanov, currents tangent to distributions, \url{https://doi.org/10.4310/jdg/1668186786}. The supplied label does not provide a full bibliographic citation.
\item Supplied literature reference: original Oberwolfach report, \url{https://doi.org/10.4171/owr/2014/33}.
\item Literature assessment dated exactly \texttt{2026-08-21}: supplied dated curation, not independently verified. It labels the topic \emph{partially solved}, describes the specific horizontal absolute-continuity question as answered positively by a later tangent-current theorem, and says that Morgan's broader normal-current decomposition and stronger distribution questions are not fully resolved. The supplied background explicitly requests independent verification of time-sensitive status and cited claims.
\item Attribution qualification: no proposer is named in the proposer field. Morgan is mentioned only in the supplied literature assessment. The registry labels the statement \texttt{corrected\_verified}, explaining that mathematical content was preserved while notation, typography, or omitted definitions were repaired; this is supplied metadata, not independent mathematical verification.
\end{itemize}
\end{document}
